<%= heading2 %>{Theorem-like Environments and Mathematics}

The \texttt{llncs} class predefines the theorem-like environments \texttt{theorem}, \texttt{proposition}, \texttt{lemma}, \texttt{corollary}, \texttt{definition}, \texttt{example}, \texttt{exercise}, \texttt{note}, \texttt{problem}, \texttt{question}, \texttt{remark}, and \texttt{solution} as well as the unnumbered \texttt{proof} environment -- no additional package is required.

<%- bexample %>
\begin{definition}[Idempotence]
  An operation $f$ is \emph{idempotent} if $f(f(x)) = f(x)$ holds for all inputs $x$.
\end{definition}

\begin{proposition}
  The identity operation is idempotent.
\end{proposition}

\begin{lemma}\label{lem:fixed}
  An idempotent operation fixes every element of its image.
\end{lemma}

\begin{theorem}\label{thm:idempotent-composition}
  If two idempotent operations $f$ and $g$ commute, their composition $f \circ g$ is idempotent.
\end{theorem}

\begin{proof}
  Since $f$ and $g$ commute and are idempotent,
  \begin{equation}
    (f \circ g)\bigl((f \circ g)(x)\bigr) = f\bigl(f(g(g(x)))\bigr) = f(g(x)) = (f \circ g)(x).
    \label{eq:idempotent-composition}
  \end{equation}
  \qed
\end{proof}

\begin{corollary}
  Applying $f \circ g$ a second time yields no further change, as \cref{eq:idempotent-composition} shows.
\end{corollary}

\begin{example}
  Normalizing whitespace in a string is idempotent: \cref{lem:fixed} applies to every already-normalized string.
\end{example}
<%- eexample %>
